\begin{abstract}
Accurate automated classification of software requirements into functional specifications and fine-grained non-functional quality attributes is fundamental to software architecture design, traceability maintenance, and compliance management. However, conventional requirement engineering pipelines remain constrained by a persistent trade-off: traditional lexical classifiers struggle with domain vocabulary sparsity, while modern deep transformer models demand substantial computational infrastructure and remain vulnerable to severe real-world class imbalance and subword truncation. Earlier bio-inspired attempts to bridge this gap relied on static, single-nucleotide mappings that failed to capture contextual syntactic dependencies, capping classification performance. To address these limitations, this study presents \textbf{DNA+}, a bio-inspired requirements engineering framework that translates natural language requirement statements into synthetic genomic sequences to exploit structural biological grammar for text categorization. The framework dynamically pairs primary requirement anchors with auxiliary syntactic dependencies, synthesizes authentic biological sequences in standard FASTA format, and extracts nucleotide transition frequencies, codon distributions, and sequence entropy. These biological attributes are fused with sentence-level semantic embeddings into a compact multimodal feature manifold and processed through an ensemble learning architecture supported by dual-tier explainable artificial intelligence (TreeSHAP and LIME) for token-level transparency. DNA+ was comprehensively evaluated across three diverse benchmark corpora encompassing 8,654 authentic software requirements (canonical PROMISE, extended P-EXP, and industrial FNFC) using a repeated stratified ten-fold cross-validation protocol across thirty randomized splits without synthetic resampling distortion. The framework achieved a peak classification accuracy of 94.77\% and a macro-averaged F1-score of 93.90\% on the industrial benchmark, outperforming baseline models by up to 17.42\%. Under severe class imbalance exceeding 70:1, DNA+ demonstrated exceptional discriminative resilience, retaining an area under the precision-recall curve of 0.9361---over eight times higher than random chance. Systematic feature ablation confirmed that biological sequence features are the primary performance driver, preventing the severe accuracy collapse observed in pure lexical models. The primary novelty of this study lies in establishing an alignment-free genomic sequence representation that captures structural requirement grammar, providing a lightweight, interpretable, and computationally efficient alternative to resource-intensive language models for software engineering.
\end{abstract}

\begin{keywords}
Requirements engineering; non-functional requirements; bio-inspired computing; text classification; genomic sequence encoding; Sentence-BERT; class imbalance; explainable artificial intelligence.
\end{keywords}
